\documentclass{jsarticle}
\usepackage{tabularx}
\usepackage{booktabs}

\begin{document}
\section*{第3部　データテーブル設計書}
\subsection*{1.　データベーステーブル}
データベースはMySQLを想定する。以下は提案書に記載された機能を実現するための論理テーブル案です。提案書では物理テーブル名、データ型、桁数、主キー、外部キーの詳細までは確定してないため、これらは内部設計で定義する。

\noindent
\begin{tabularx}{\textwidth}{|l|X|X|X|}
    \hline
    \textbf{No.} & \textbf{テーブル名} & \textbf{主な項目（案）} & \textbf{用途} \\ 
    \hline
    1 & 利用者テーブル & 利用者ID、メールアドレス、パスワードハッシュ、利用者種別、登録日時 & 主催者・競技者の認証と権限管理 \\ \hline
    2 & 競技者テーブル & 競技者ID、利用者ID、プレイヤー名、公開プロフィール情報 & 競技者を固定IDで識別し、各局実践後も履歴を関連付ける \\ \hline
    3 & 大会テーブル & 大会ID、主催者ID、大会名、対戦ゲーム、開催日時、開催場所、最大参加人数、大会形式、受付期間、開催状態、配信URL & 大会基本情報の管理 \\ \hline
    4 & 大会参加テーブル & 参加ID、大会ID、選手ID、登録日時、登録状態、受付備考 & 大会への参加登録と受付状況の管理 \\ \hline
    5 & トーナメントテーブル & トーナメントID、大会ID、大会形式、作成日時 & 大会ごとのトーナメント情報の管理 \\ \hline
    6 & 試合テーブル & 試合ID、トーナメントID、ラウンド、試合番号、試合状態、予定日時 & 対戦カードと試合進行状況の管理 \\ \hline
    7 & 試合競技者テーブル & 試合ID、競技者ID、対戦枠、得点、勝敗 & 試合と競技者の対応付け \\ \hline
    8 & 試合結果報告テーブル & 報告ID、試合ID、入力者ID、入力特異、確認状態、確認日時 & 競技者同士の結果報告と主催者確認の管理 \\ \hline
    9 & 競技者お気に入りテーブル & 利用者ID、競技者ID、登録日時 & 観戦者のお気に入り競技者の管理 \\ \hline
    10 & 通知履歴テーブル & 通知ID、利用者ID、通知種別、通知内容、送信日時、送信状態 & メール通知の送信履歴管理 \\ \hline
    11 & 対戦履歴情報 & 試合ID、競技者ID、対戦相手ID、結果 & 選手間の過去対戦情報の検索（試合情報から導出する方式も可） \\ \hline
    
\end{tabularx}

\vspace{1em}
\noindent
\textbf{注：}上記の項目は機能要件から整理した論理設計案であり、確定したデータベース定義ではない。特に、観戦者アカウントを必須とするか、試合結果の入力・確認履歴をどの粒度で保持するか、対戦履歴を独立テーブルにするかは内部設計で決定する。
\end{document}